\section{R instruments and credit thresholds: preserving the design and its limits}
\label{exec22:riv-scope}

A credit eligibility threshold can change access to a product without forcing
every eligible customer to take it. A bank comparing losses above and below the
threshold must decide whether its question concerns eligibility, actual take-up,
the customers whose take-up changes, or universal provision of the product. These
are different interventions and populations. An R interface that returns one
coefficient under the label ``causal effect'' can conceal the most consequential
choice in the analysis. The useful comparison is therefore between complete
designs: the assumptions, target, statistical procedure and information that must
survive a library, CLI or MCP request.

Six concrete package profiles supply complementary pieces. AER~1.2-17 and
ivreg~0.6-8 provide separately implemented instrumental-variable regression
interfaces. ivmodel~1.9.1 adds weak-instrument tests and sensitivity procedures
for a specified linear model. rdrobust~4.0.0 supplies local-polynomial estimation
and bias-corrected inference at a cutoff; rddensity~3.0 supplies density diagnostics;
rdlocrand~3.0 supplies methods under local randomization. AER remains relevant
both as a comparison and as a declared dependency of rdlocrand. The shared
\texttt{ivreg} function and class names do not make the first two packages the
same engine.~\citep{exec22_riv_aer,exec22_riv_ivreg,exec22_riv_ivmodel,
exec22_riv_rdrobust,exec22_riv_rddensity,exec22_riv_rdlocrand}

The release date also matters. rdrobust~4.1.1 carries the package date
7~October~2026, but its CRAN publication field is 23:20:02~UTC on that date,
after the client cutoff of 15:59:59~UTC. The cutoff comparison therefore retains
the earlier authenticated 4.0.0 distribution. This resolves the CRAN distribution
question without asserting that every possible upstream channel had the same
publication history. Current manuals cannot silently supply capabilities to an
earlier version.~\citep{exec22_riv_rdrobust}

\subsection{What the instrument identifies, and which policy it values}
\label{exec22:riv-target}

Let $Z\in\{0,1\}$ denote eligibility or a randomized encouragement, $D(z)$
actual take-up under eligibility $z$, and $Y(d)$ a fixed-horizon monetary loss
under treatment $d$. Loss is positive when it costs the bank money. In an
independent-encouragement design, assume consistency, no interference,
independence of $Z$ from the potential variables, exclusion of a direct $Z$
effect on loss, monotonicity $D(1)\geq D(0)$, and a positive first stage.
Write $C=\{D(1)>D(0)\}$ for the complier group. Independence and exclusion give
\begin{align}
 \Delta_D&=E[D\mid Z=1]-E[D\mid Z=0]=P(C),\nonumber\\
 \Delta_Y&=E[Y\mid Z=1]-E[Y\mid Z=0]
          =P(C)E[Y(1)-Y(0)\mid C],\nonumber\\
 \theta_C&=\Delta_Y/\Delta_D.
 \label{exec22:riv-late}
\end{align}
The second equality follows because always-takers and never-takers do not change
treatment, while every complier changes from zero to one. It does not identify
their unobserved treatment contrasts, and therefore does not generally identify
the population average effect. With a deterministic score cutoff, global
independence of eligibility is unavailable: an appropriate continuity argument
at the cutoff, or an explicitly justified local-randomization window, must supply
the corresponding design. Calling eligibility an instrument does not supply that
argument.

Consider the following complete, project-authored population. The encouragement
is independent of type and has probability $1/2$. All monetary values are
hypothetical dollars over the same horizon.~\citep{exec22_riv_reference}
\begin{equation}
\begin{array}{c|c|cc|rr}
 \text{type}&\text{mass}&D(0)&D(1)&Y(0)&Y(1)\\\hline
 \text{complier}&2/5&0&1&200&50\\
 \text{always-taker}&1/5&1&1&200&150\\
 \text{never-taker}&2/5&0&0&200&200
\end{array}
\label{exec22:riv-population}
\end{equation}
Mean take-up is $1/5$ without encouragement and $3/5$ with it. Mean loss changes
from USD~190 to USD~130. Thus $\Delta_D=2/5$, $\Delta_Y=-60$ and
$\theta_C=-150$, while the full population effect is $-70$. If treatment costs
USD~100, the stipulated complier treatment margin is USD~50, but treating everyone
instead of no one has margin USD~$-30$. Complier membership is a latent response
type, so the first margin does not imply an implementable policy that identifies
and treats only compliers.

The encouragement itself is an implementable alternative under this law. If it
costs USD~10 per customer and changes treatment costs only through the extra
take-up, its incremental margin relative to no encouragement is
\begin{equation}
 V_Z=-\Delta_Y-100\Delta_D-10=10\quad\text{USD}.
 \label{exec22:riv-value}
\end{equation}
This calculation explains why the proposed service should return the reduced
form and first stage together with their ratio. A business request about changing
eligibility may be answered by the eligibility effect without extrapolating a
complier effect to universal treatment. In actual estimates, uncertainty in both
components and their covariance must enter the policy calculation.

Two violations show what software cannot repair by fitting more accurately.
If encouragement directly lowers every customer's loss by USD~20, holding the
treatment law fixed, the observed ratio becomes $(-60-20)/(2/5)=-200$ although
the treatment contrast among compliers remains $-150$. Alternatively, let
compliers have mass $2/5$ and effect $-150$, defiers have mass $1/5$ and effect
$-50$, and the remaining customers never take treatment. Independence and
exclusion then imply a first stage $1/5$, reduced form $-50$ and ratio $-250$.
The ratio is outside the two changing groups' effects because the defier weight
is negative. Neither an overidentification test nor a first-stage statistic
establishes the exclusion and monotonicity assumptions used in
equation~\eqref{exec22:riv-late}. These finite examples isolate the failure;
they are not empirical evidence of either violation at a bank.

\subsection{Regression interfaces and weak-instrument sets}
\label{exec22:riv-inference}

AER documents a two-part regressor/instrument formula, with exogenous regressors
also included as instruments. Its summary method defaults to diagnostics being
off; when requested it reports first-stage, endogeneity and overidentification
diagnostics. The dedicated ivreg package offers additional formula conventions,
including separate exogenous, endogenous and excluded-instrument parts, and
ordinary two-stage least squares or M/MM robust regression. Its summary defaults
to diagnostic reporting, and an alternative covariance must be supplied as a
function when auxiliary diagnostic regressions also require it. A saved matrix
for the final regression does not specify covariance for those auxiliary fits.
\citep{exec22_riv_aer,exec22_riv_ivreg}

The integration should therefore name the package and version, construct an
allowed formula from typed variable roles, and record the realized design
matrices, row mask, contrasts, weights, intercept, numerical ranks and covariance
method. It should retain aliased or unavailable coefficients with their reasons.
Automatically dropping an absorbed or collinear treatment and reporting zero
would change an unavailable target into a false scientific result. The M/MM
options address outlier-resistant regression; heteroskedastic or cluster covariance
addresses other uncertainty issues. Neither label by itself supplies inference
valid under a weak first stage.

ivmodel describes a continuous-outcome model with one endogenous variable,
$Y=X\alpha+D\beta+\epsilon$, and $E[\epsilon\mid X,Z]=0$. Some procedures,
including its conditional likelihood ratio path, add a linear first-stage model.
The Anderson--Rubin and conditional likelihood ratio interfaces return a matrix
whose rows are interval components, together with a human-readable description.
The pinned source explicitly constructs whole-line and disconnected outputs.
Those are substantive answers, not serialization defects.
\citep{exec22_riv_ivmodel}

An exact Gaussian reference makes the required representation transparent. Let
$(\widehat\Delta_Y,\widehat\Delta_D)$ be jointly normal with mean
$(\beta\pi,\pi)$ and known positive-definite covariance
\begin{equation}
 V=\begin{pmatrix}v_Y&c\\c&v_D\end{pmatrix}.
 \label{exec22:riv-gaussian}
\end{equation}
For a candidate $b$, the standardized contrast
$(\widehat\Delta_Y-b\widehat\Delta_D)/\sqrt{v_Y-2bc+b^2v_D}$ is standard
normal under $\beta=b$, including when $\pi=0$. Inverting a squared critical
value $q$ therefore gives
\begin{equation}
 \mathcal C_q=\left\{b:
 (\widehat\Delta_D^2-qv_D)b^2+
 (-2\widehat\Delta_Y\widehat\Delta_D+2qc)b+
 (\widehat\Delta_Y^2-qv_Y)\leq0\right\}.
 \label{exec22:riv-set}
\end{equation}
This is a project derivation for a known-covariance Gaussian law. It is neither
a reproduction of the package's finite-sample implementation nor a guarantee
under arbitrary heteroskedasticity or clustering. Choose $q=4$ for exact simple
roots: the model coverage is $2\Phi(2)-1\simeq0.9545$, not exactly 95\%.
With $c=0$, the following observed inputs produce three different shapes:
\begin{equation}
\begin{array}{cc|cc|l}
 \widehat\Delta_Y&\widehat\Delta_D&v_Y&v_D&\mathcal C_4\\\hline
 2&0&3/4&1/4&(-\infty,-1]\cup[1,\infty)\\
 2&2&5/12&1/4&[1/3,7/3]\\
 0&0&3/4&1/4&\mathbb R
\end{array}
\label{exec22:riv-topology}
\end{equation}
Direct substitution and quadratic roots agree in each case.
\citep{exec22_riv_reference}

For the first row, searching only $[-1/2,1/2]$ finds no accepted value even
though the true inverted set contains two infinite rays. An interface must
distinguish ``empty on the searched grid'' from a globally empty set. It must
also distinguish an infinite endpoint from missing information and from an
unfinished numerical search. Replacing the two components by their convex hull
would insert all the rejected values between $-1$ and $1$; replacing them by a
finite interval would remove supported tail values. A monetary decision can
change under either alteration.

Static inspection adds a specific implementation concern. In the pinned
ivmodel CLR routine, a numerical error during conditional p-value evaluation is
caught and replaced with zero; a root-search error similarly substitutes zero
for an intermediate quantity. The AR routine uses a projection F statistic,
while CLR constructs a residual cross-product covariance. The parent object's
generic heteroskedastic or cluster standard-error options do not establish that
these paths implement the corresponding robust tests.~\citep{exec22_riv_ivmodel}
A proposed worker must expose or independently detect the numerical failure
before accepting the result. An explicit zero p-value is otherwise
indistinguishable from the substituted failure value at that field alone. The
source establishes the branch; its frequency and native triggering cases remain
to be measured in a controlled prototype. This is a reason to retain the
package as a reference while withholding runtime admission for that path.

\subsection{A local polynomial cannot fill a gap in the score support}
\label{exec22:riv-continuity}

rdrobust's continuity-based profile distinguishes sharp, fuzzy and kink designs,
the degree used for estimation and bias correction, left/right estimation
bandwidths $h$, bias bandwidths $b$, kernel, variance procedure and clustering.
Its outputs distinguish conventional estimates, bias-corrected estimates and
the associated uncertainty rows. The adapter must preserve these associations;
a corrected estimate with the wrong standard error is a different procedure.
Effective sample sizes, unique score counts and warnings are part of the result.
\citep{exec22_riv_rdrobust}

Rounded credit scores make the support issue concrete. Suppose the only score
values are $R=-1$ and $R=1$, eligibility is $D=1\{R\geq0\}$, and observed mean
losses are USD~100 and USD~80 respectively. Two continuous potential-mean models
agree at every observed score:
\begin{align}
 \text{A: }&m_0(r)=100,\qquad m_1(r)=80,\nonumber\\
 \text{B: }&m_0(r)=m_1(r)=90-10r.
 \label{exec22:riv-support}
\end{align}
Model A has treatment effect USD~$-20$ at zero, while model B has effect zero.
There are no observations approaching the cutoff from either side, so continuity
alone cannot select between the models. More customers at the same two scores
reduce uncertainty about the observed means but do not identify the unobserved
limits. A mass-point option can alter estimation behavior; it cannot create
missing score support.~\citep{exec22_riv_reference,exec22_riv_rdrobust}

If additional substantive information bounds each potential-mean slope in
absolute value by $L$, and the nearest observations are distance $\delta$ from
the cutoff, each extrapolation contributes at most $L\delta$. In this
two-point law the cutoff effect belongs to
\begin{equation}
 \theta(0)\in[-20-2L\delta,-20+2L\delta].
 \label{exec22:riv-bound}
\end{equation}
The endpoints are attainable by choosing the two potential-mean slopes within
the declared bound. With $L=10$ and $\delta=1$, the effect set is
$[-40,0]$ dollars. At an assumed treatment cost USD~10, the net-margin set is
$[-10,30]$ dollars, and it does not justify a uniformly positive-margin claim.
These are identification bounds with known observed means; sampling uncertainty
would need to be added rather than silently absorbed into the slope restriction.

rddensity supplies a different diagnostic: local-polynomial density estimates
and tests for a discontinuity in the running-variable distribution, with explicit
bandwidth, variance and mass-point options. Its documentation distinguishes a
conventional statistic from robust bias correction and warns about conventional
inference with an MSE-selected bandwidth.~\citep{exec22_riv_rddensity}
Even exact equality of the score density on both sides would not prove continuity
of potential losses. For example, a continuously uniform score can coexist with
another policy that changes at zero and affects loss, even when the treatment
being studied has no effect. A density test examines the score law, not that
competing intervention. Conversely, rounded scores need their own support and
measurement discussion before a continuous-density interpretation is applied.

\subsection{Local randomization is an alternative design with its own obligations}
\label{exec22:riv-randomization}

rdlocrand assumes local randomization within a supplied or selected window. The
window-selection routine considers nested windows and covariate balance tests.
That procedure is useful evidence for design assessment, but a favorable balance
p-value cannot prove assignment independence. Outcome-driven window search also
changes the inference problem. The proposed request must retain whether the window
was fixed before outcome inspection, selected with pre-treatment covariates, or
chosen through a procedure requiring additional selection analysis.
\citep{exec22_riv_rdlocrand}

The distinction is visible in a finite reference. Fix six customers' losses under
the sharp null at $(0,0,0,100,100,100)$ dollars. Assume exactly three are assigned
treatment, uniformly over the $\binom63=20$ allocations, in a window fixed in
advance. If the first three receive treatment, the difference in mean loss is
$-100$. Exactly two allocations have an absolute difference at least this large,
so the two-sided randomization p-value is $2/20=0.1$. This is complete enumeration
under a Fisher sharp null of no effect for any customer. It does not supply a
general test of zero population average effect under arbitrary heterogeneous
effects.~\citep{exec22_riv_reference}

The pinned rdlocrand documentation also distinguishes requested from effective
polynomial degree. For the documented mean contrasts and fuzzy procedures,
positive degree invokes large-sample normal inference; a randomization p-value
is not computed. Some statistic/degree combinations revert to degree zero with
a warning. Fuzzy test inversion can require an explicit effect grid. Consequently,
the adapter must export the effective method, degree, variance choice, assignment
law, seed, repetitions, window and grid, not just the originally requested
options. HC3 standard errors do not confer finite-sample randomization validity,
and a TSLS/Wald label does not remove weak-first-stage concerns.
\citep{exec22_riv_rdlocrand}

\subsection{Reuse choices, rights and a prototype that can change the decision}
\label{exec22:riv-reuse}

The immediate recommendation is reference reuse with conditional plans for
original R workers. The dedicated ivreg interface is the natural initial
ordinary-2SLS candidate; AER remains a distinct baseline and dependency.
ivmodel is valuable for documenting sets and sensitivity, with CLR admission
conditional on a qualified failure path. rdrobust serves an explicit continuity
profile, rddensity a supplementary diagnostic, and rdlocrand a separate
local-randomization profile. A direct port adds numerical and maintenance duties
and should follow an identified operational need. Existing grf effect-estimation
and fixest panel-worker proposals do not automatically qualify an IV target or
this threshold design merely because they share an R execution path.

The pinned metadata declares GPL version alternatives for AER, GPL version~2 or
later for ivreg, and GPL~3 for the three RD packages. ivmodel declares GPL~2 or
its local license file, whose inspected text supplies an MIT alternative with
copyright/permission-notice conditions and a warranty disclaimer. These facts
improve the package-level comparison but do not clear the full dependency tree.
rdlocrand imports AER and sandwich; other profiles also depend on statistical,
formula, matrix or plotting packages. Exact installed versions and their notices
must be included in the selected distribution.~\citep{exec22_riv_aer,
exec22_riv_ivreg,exec22_riv_ivmodel,exec22_riv_rdrobust,exec22_riv_rddensity,
exec22_riv_rdlocrand}

The GPL texts distinguish running a program from copying, modifying and conveying
it, and impose applicable notice and corresponding-source conditions. Moving an
engine behind a process boundary does not itself decide all obligations of a
particular distributed product. Original-worker integration and porting therefore
need a review of the actual distribution, including third-party material. No
current package recommendation grants that implementation clearance.
\citep{exec22_riv_gpl2,exec22_riv_gpl3}

A useful public prototype would first compare typed requests with independent
projection and set-inversion references, then test the complete native workflow.
The ordinary-IV cases should vary rank, weights, missingness and covariance
conventions while preserving the target. Weak-IV cases should stress a vanishing
first stage, disconnected sets, numerical integration and root-search failures.
Threshold cases should distinguish smooth support approaching zero, a small
take-up discontinuity, rounded scores and competing policies. The local-randomization
path should reproduce the twenty-allocation example before larger simulations,
and analyze selected-window behavior separately.

Before final simulation, freeze the data laws, target, package/dependency versions,
seeds, replication budget, tolerances and stopping rule. Keep failed fits in the
denominator and report Monte Carlo uncertainty. Native execution, serialization
fidelity, statistical calibration and decision value answer different questions;
passing one must not silently mark the others complete. An LLM can help translate
the business question, propose alternative eligibility graphs and assemble these
typed requests. It must preserve unresolved exclusion, compliance and score-support
questions, and must not convert a failed search or a diagnostic p-value into a
causal certificate. This design makes a bounded prototype informative before
an expensive experiment with bank data.
